\chapter{Introduction}\label{chap:intro}

Consider a cybersecurity student who has just completed a network penetration testing course and wants to explore hardware security. She knows that real-world attacks against IoT devices---extracting firmware through debug interfaces, identifying hardcoded credentials on flash storage, exploiting bootloader misconfigurations---require hands-on interaction with physical circuit boards. Yet the path forward is daunting. A basic workstation with a multimeter, logic analyser, UART adapter, and SPI programmer costs several hundred euros. Expendable target devices must be sourced and periodically replaced. Professional training courses at conferences like Black Hat or through SANS carry fees between 3{,}000 and 8{,}000~EUR per participant~\cite{genova2025}. Unless her university maintains a dedicated hardware security laboratory---and most do not---she has no affordable way to develop these skills.

This scenario is not hypothetical---it captures a structural problem in cybersecurity education with measurable consequences. The Mirai botnet of 2016~\cite{antonakakis2017mirai} exploited default credentials on hundreds of thousands of consumer IoT devices, causing some of the largest distributed denial-of-service attacks in internet history. Costin et al.~\cite{costin2014large} analysed over 32{,}000 firmware images from embedded devices and found that a majority contained at least one serious vulnerability---hardcoded passwords, unencrypted storage, or exposed debug interfaces. These are not exotic attack surfaces requiring nation-state resources; they are entry-level findings that a trained analyst can identify with a multimeter, a serial adapter, and basic command-line skills. The bottleneck is not the difficulty of the analysis but the scarcity of training environments in which students can develop these skills without access to a physical laboratory. On the software side, a mature ecosystem of Capture The Flag platforms has emerged: CTFd~\cite{ctfd}, PicoCTF~\cite{picoctf}, Hack The Box~\cite{hackthebox}, and TryHackMe~\cite{tryhackme} deliver scalable, browser-accessible challenges covering web exploitation, binary analysis, cryptography, and reverse engineering. Hamari et al.\ showed that gamification elements---points, leaderboards, progressive unlocks---significantly increase engagement in learning contexts~\cite{hamari2014}, and these platforms exploit that insight to great effect. Svabensky et al.'s comprehensive survey of hands-on cybersecurity training confirms that practical, interactive exercises produce measurably better learning outcomes than lecture-only instruction~\cite{svabensky2020}. Yet their survey also reveals a stark asymmetry: virtually all existing platforms target software vulnerabilities. Hardware security training remains locked behind equipment costs, logistics, and physical access requirements.

Intermediate solutions exist but leave significant gaps. Remote laboratories such as iLab~\cite{ilab2004} and WebLab-Deusto~\cite{weblab2005} grant network access to real instruments, yet they demand dedicated server infrastructure and support only limited concurrency. Hardware CTF competitions like RHme~\cite{rhme} and the Embedded Security Challenge at CSAW~\cite{csawesc} produce excellent pedagogical experiences, but they are episodic events that require physical kits shipped to registered participants. Platforms like Microcorruption~\cite{microcorruption} and OWASP IoTGoat~\cite{iotgoat} focus narrowly---the former on a single microcontroller architecture, the latter on network-layer services---and neither simulates the physical board-level workflow that defines hardware analysis in practice: visually inspecting a PCB, probing test points with a multimeter, wiring a serial adapter with the correct crossover convention, dumping firmware through an SPI programmer. No existing solution combines these capabilities with a no-code authoring workflow in a zero-cost, browser-based package.

This thesis addresses that gap. We present \pcbctf{}, a web-based interactive cyber range for hardware security training in CTF format. The platform enables a student to interact with a high-resolution photograph of a real printed circuit board, identify components through visual inspection and magnification, perform simulated electrical measurements with a digital multimeter, establish serial connections through a virtual UART adapter, extract firmware via a simulated SPI programmer, and explore a fully emulated embedded terminal---all within the browser, with no server-side code execution.

The first contribution is an architecture for hardware instrument simulation that runs entirely on the client side. The simulated multimeter, UART adapter, and SPI programmer produce realistic electrical feedback---voltage levels, impedance values, protocol-level validation of crossover wiring and role-based probe matching---while a normalised coordinate system ensures that interactive overlays scale responsively across display sizes. Because all simulation logic executes in the browser, the platform eliminates server-side security concerns and operates offline after the initial page load.

The second contribution is a configurable terminal emulation system. We design a declarative configuration language that allows exercise authors to specify virtual filesystems, multi-stage boot sequences (including U-Boot and Linux login stages), custom commands with multiple output types, constraint-based access control, and a progressive flag discovery mechanism. The terminal reproduces the experience of interacting with a real embedded device through a synchronous command pipeline, without running any actual shell.

The third contribution is a no-code authoring interface. A visual exercise editor allows instructors to configure complete training scenarios---PCB image markup, component and pin placement, step and objective definition, terminal environment setup, and preset management---through a graphical interface. This lowers the barrier for content creation and enables educators without web development experience to produce and share exercises.

The fourth contribution is a validation through a real-world case study. We construct a complete exercise based on the hardware analysis of a TP-Link WR841N router, building directly on the vulnerability assessment methodology and device teardown conducted by Genova~\cite{genova2025}. The exercise demonstrates that a multi-step investigation---from physical reconnaissance through UART connection to vulnerability discovery---can be faithfully reproduced in a fully virtual environment.

This work has been conducted within the \artic{} project, part of Spoke~4 (``Operating Systems and Virtualization Security'') of the SERICS (Security and Rights in CyberSpace) national research programme, funded under the Italian PNRR initiative.

Chapter~\ref{chap:background} establishes the technical foundations of hardware security, embedded systems, and common debug interfaces. Chapter~\ref{chap:stateoftheart} surveys the current landscape of hardware security training, CTF platforms, and remote laboratories. Chapter~\ref{chap:contributions} formalises the requirements that drive the platform design, while Chapter~\ref{chap:architecture} details the system architecture and Chapter~\ref{chap:implementation} covers implementation. Chapter~\ref{chap:casestudy} presents the TP-Link WR841N case study, and Chapter~\ref{chap:conclusions} draws conclusions and identifies directions for future work. Taken together, these chapters argue that accessible, high-fidelity hardware security training can be delivered through a browser-based platform at zero marginal cost per student---and that doing so is not merely convenient but pedagogically necessary.
